\documentclass[11pt,a4paper]{article}

\usepackage[utf8]{inputenc}
\usepackage[T1]{fontenc}
\usepackage{lmodern}
\usepackage[margin=1in]{geometry}
\usepackage{amsmath,amssymb,amsthm}
\usepackage{booktabs}
\usepackage{array}
\usepackage{verbatim}
\usepackage{hyperref}
\usepackage{enumitem}
\usepackage{microtype}

\hypersetup{
  colorlinks=true,
  linkcolor=black,
  urlcolor=blue,
  citecolor=blue,
  pdftitle={SiftTSP: A Geometric Sandclock-Sift Heuristic for Euclidean TSP}
}

\newcommand{\bigO}{\mathcal{O}}

\title{SiftTSP: A Geometric Sandclock-Sift Heuristic\\ for Euclidean TSP}
\author{Norbert Nopper}
\date{}

\begin{document}
\maketitle

\begin{abstract}
We present \textbf{SiftTSP}, a deterministic polynomial-time heuristic for the
Euclidean Travelling Salesman Problem in the partitioning lineage of
Karp~\cite{karp1977}. The algorithm partitions the city set by recursive geometric
bisection into $2^d$ sections, solves small sections by brute force and
larger ones by a complementary pair of greedy rules (min-max / max-min),
and reconnects the section paths by exhaustive search over orderings and
orientations. The angular dimension of the search is covered by two
complementary sweeps --- a uniform \emph{mill} (breadth) and a damped
fan-shaped \emph{sieve} (depth) --- coordinated by a \emph{sandclock}
depth schedule $1 \to d_{\max} \to 1$ and an 8-state Phase~4 loop over
(section-mode, mill-direction, ordering). The contribution is the
specific orchestration of these primitives (\S1.5), not a new
theoretical guarantee.

For any fixed parameters the algorithm runs in $O(n^2)$ time
(Theorem~1). SiftTSP is \textbf{not exact}: we exhibit an explicit
12-point Euclidean instance (\texttt{counterexample.py}, \texttt{rings-12})
on which the algorithm misses the Held--Karp optimum at every tested
depth, with gap $+3.24\%$ at $d=3$. The structural reason is identified
as \emph{decomposition irrecoverability} (\S2.3): no fixed family of
geometric decompositions can recover an optimal tour that interleaves
cities across components of every decomposition in the family.
Consequently this work makes \textbf{no claim} of resolving
$\mathcal{P}$ versus $\mathcal{NP}$~\cite{cook1971,levin1973,aaronson2009}, and offers no
approximation ratio (unlike Christofides~\cite{christofides1976} for metric TSP or
Arora~\cite{arora1998} for Euclidean TSP).

On a battery of 14 instances ($n \in \{12, 14, 16\}$) including zigzag,
comb, interleaved, concentric-ring, and uniform configurations, SiftTSP
matches Held--Karp on most structured inputs and beats 2-opt on the
rotationally symmetric \texttt{rings-12} and \texttt{rings-14} instances
despite not being exact.

\medskip
\noindent\textbf{Keywords:} Travelling Salesman Problem, Polynomial-Time
Heuristic, Euclidean TSP, Geometric Subdivision, Mill-Sieve Search.
\end{abstract}

\section{Algorithm}

\subsection{Setting}

Given $n$ cities $C = \{c_1, \dots, c_n\} \subset \mathbb{R}^2$, a
\emph{tour} is a cyclic permutation $\pi$ of $C$, and its length is
$L(\pi) = \sum_i \Vert c_{\pi(i)} - c_{\pi(i+1)} \Vert_2$ (indices
$\bmod\ n$). Euclidean TSP asks for an $L$-minimum tour. The decision
form of general (metric) TSP is NP-Complete by Karp's reduction from
Hamiltonian Cycle~\cite{karp1972}; Papadimitriou~\cite{papadimitriou1977}
established that the restriction to Euclidean instances in the plane is
also NP-Complete, hence the optimisation form is NP-Hard. The exact
Held--Karp DP~\cite{heldkarp1962} runs in $O(n^2 \cdot 2^n)$ time and is
feasible only for small $n$. SiftTSP is a polynomial-time
\emph{heuristic} --- it is not exact, see \S2.3.

\subsection{Parameters and Phases}

SiftTSP is controlled by four constants treated as fixed:

\begin{center}
\begin{tabular}{ll}
\toprule
Symbol & Meaning \\
\midrule
$d_{\text{ceiling}}$ & Maximum recursive subdivision depth \\
$m_{\text{ceiling}}$ & Mill (breadth) cap: third-steps per direction \\
$s_{\text{ceiling}}$ & Sieve (depth) cap: damped halvings around mill best \\
$\tau$ & Brute-force threshold: sections with $\leq \tau$ cities are solved exactly \\
\bottomrule
\end{tabular}
\end{center}

The algorithm is structured in four phases:

\begin{itemize}
\item \textbf{Phase~1 (max).} A worst-case greedy-longest tour, computed
once, fixes an initial upper bound $B_0 \geq L^{\star}$ that the rest only
lowers.

\item \textbf{Phase~2 (sift).} At each depth $d$ on the sandclock
schedule $1, 2, \dots, d_{\max}, \dots, 2, 1$ and at each angle $\theta$
in the current mill+sieve schedule, the city set is partitioned by
recursive median bisection along \textbf{each} of four axis-pairs (each
pair: two perpendicular axes from the rotated frame at angle $\theta$
--- the primary axis governs the top-level split, the secondary axis the
next level, alternating thereafter). The four pairs are taken from
$\{\text{vertical},\ /,\ \text{horizontal},\ \backslash\}$. Each
axis-pair yields a partition into $2^d$ sections, evaluated
independently. Sections of size $> \tau$ are solved by a greedy section
solver in one of two complementary modes (selected per sweep by
Phase~4):
\begin{itemize}
\item \textbf{min-max}: at each step, append the city minimising the
running \emph{maximum} edge (favouring smooth paths);
\item \textbf{max-min}: at each step, append the city maximising the
running \emph{minimum} edge (a complementary ordering).
\end{itemize}

\item \textbf{Phase~3 (min).} Sections of size $\leq \tau$ are solved
exactly by brute force. The $2^d$ section paths are reconnected by
exhaustive search over $(2^d - 1)! \cdot 2^{2^d - 1}$ orderings and
orientations.

\item \textbf{Phase~4 (loop).} An 8-state super-cycle drives the angular
search.
\end{itemize}

Both greedy rules are seeded at the first city in the section's input
ordering and extend the path one vertex at a time; consequently SiftTSP
is deterministic but its output depends on the order in which the
cities are presented (the partitioning, the section orderings inside
each partition, and the brute-force fall-backs all inherit this
order-dependence).

The angular dimension uses two complementary sweeps. The \textbf{mill}
is a uniform third-step rotation across $[0, \pi/2)$ in both directions
(breadth). The \textbf{sieve} is a damped forth-and-back oscillation
around the best $(\theta^{\star}, d^{\star})$ found by the mill, with
amplitudes $\tfrac{s_0}{2}, \tfrac{s_0}{4}, \dots, \tfrac{s_0}{2^{s}}$
where $s_0 = (\pi / 2^{d^{\star}}) / 3$ (depth). Mill and sieve refine
the \emph{same} angular axis at different scales; together they sift
the search through both wide and narrow neighbourhoods of every depth.

\subsection{The Phase~4 Super-Cycle}

Each sandclock sweep is parameterised by a tuple
$(d_{\max}, m, s, \mu, \delta)$, where $\mu$ is the section-solver mode
(min-max / max-min) and $\delta$ is the angular-sweep direction
(CW / CCW). The four configurations $(\mu, \delta)$ may be visited in
two orderings --- \emph{mode-first} (MD) or \emph{direction-first} (DM)
--- giving an 8-state super-cycle:

\begin{center}
\begin{tabular}{cccc}
\toprule
$k$ & $\mu$ & $\delta$ & ordering \\
\midrule
0 & max-min & CW  & MD \\
1 & min-max & CW  & MD \\
2 & max-min & CCW & MD \\
3 & min-max & CCW & MD \\
4 & max-min & CW  & DM \\
5 & max-min & CCW & DM \\
6 & min-max & CW  & DM \\
7 & min-max & CCW & DM \\
\bottomrule
\end{tabular}
\end{center}

The order of visitation matters because the mill cap $m$ and sieve cap
$s$ double \emph{only on improvement}; the sequence of $(\text{config},
m, s)$ triples ever evaluated depends on the order.

\paragraph{Transition rule.} After a sweep at position $k$:
\begin{enumerate}
\item If the \emph{mill} improved the running best, $m \leftarrow 2m$.
\item If the \emph{sieve} improved the running best, $s \leftarrow 2s$.
\item Advance $k \leftarrow (k+1) \bmod 8$.
\item If 8 consecutive sweeps at the same $(d_{\max}, m, s)$ all fail,
bump $d_{\max} \leftarrow d_{\max} + 1$, reset $m, s, k$ to $1, 1, 0$,
and restart Phase~2.
\end{enumerate}

\paragraph{Termination.} For fixed parameters
$(d_{\text{ceiling}}, m_{\text{ceiling}}, s_{\text{ceiling}})$, the
candidate tours examined at any sweep are determined by the discrete
configuration $(d_{\max}, m_{\text{eff}}, s_{\text{eff}}, \mu, \delta,
k \bmod 8)$, drawn from an $O(1)$-sized set whose size depends only
on the parameters and not on $n$. The running best $L$ is updated
only on \emph{strict} improvement, so each improving sweep contributes
a previously unseen, strictly shorter candidate; the total number of
improving sweeps is therefore bounded by the size of this candidate
set, hence $O(1)$. Each non-improving sweep increments
\texttt{no\_improve}; reaching $8$ breaks the inner \texttt{while}.
Between any two consecutive improvements there are at most $7$
non-improving sweeps, and at most $8$ non-improving sweeps follow the
last improvement, so the total number of sweeps per $d_{\max}$ is
also $O(1)$. The outer \texttt{for} loop runs over $d_{\max} = 1,
\dots, d_{\text{ceiling}}$, so the algorithm terminates after $O(1)$
sweeps overall (for fixed parameters). The returned tour has length
monotone non-increasing along the execution trace and is bounded
below by $0$.

\subsection{Pseudocode}

\begin{small}
\begin{verbatim}
SUPERCYCLE := concat(
    [(MaxMin,CW), (MinMax,CW), (MaxMin,CCW), (MinMax,CCW)],   # ordering MD
    [(MaxMin,CW), (MaxMin,CCW), (MinMax,CW), (MinMax,CCW)]    # ordering DM
)                                                              # length 8

procedure SiftTSP(C, d_ceiling, m_ceiling, s_ceiling, tau):
    best <- WorstCaseTour(C); best_len <- Length(best)         # PHASE 1
    for d_max = 1 .. d_ceiling do                              # PHASE 4
        m, s, k, no_improve <- 1, 1, 0, 0
        while m <= m_ceiling or s <= s_ceiling do
            (mode, dir) <- SUPERCYCLE[k]
            (best, best_len, mill_imp, sieve_imp)
              <- SandclockSweep(C, d_max, min(m,m_ceiling),
                                min(s,s_ceiling), tau, mode, dir,
                                best, best_len)
            k <- (k+1) mod 8
            if mill_imp or sieve_imp then
                no_improve <- 0
                if mill_imp  and m <= m_ceiling then m <- 2*m
                if sieve_imp and s <= s_ceiling then s <- 2*s
            else
                no_improve <- no_improve + 1
                if no_improve >= 8 then break
    return best

procedure SandclockSweep(C, d_max, m, s, tau, mode, dir, best, best_len):
    mill_imp, sieve_imp, theta*, d* <- false, false, 0, d_max
    for d in [1,2,..,d_max,..,2,1] do                          # PHASE 2 mill
        for theta in MillAngles(d, m, dir) do
            (best, best_len, hit) <- EvalAt(C, d, theta, tau, mode, ...)
            if hit then mill_imp, theta*, d* <- true, theta, d
    s0 <- (pi / 2^d*) / 3                                      # PHASE 2 sieve
    for j = 0 .. s-1 do
        delta <- s0 / 2^(j+1)
        for sigma in (+1, -1) do
            (best, best_len, hit) <- EvalAt(C, d*, theta* + sigma*delta, ...)
            if hit then sieve_imp, theta* <- true, theta* + sigma*delta
    return (best, best_len, mill_imp, sieve_imp)

procedure EvalAt(C, d, theta, tau, mode, best, best_len):
    hit <- false
    for (axis_p, axis_s) in AxisPairs(theta) do                # 4 pairs
        sections <- Subdivide(C, [axis_p, axis_s], d)
        paths <- []
        for S in sections do
            if |S| <= tau then paths <- paths + [BruteForce(S)]      # PHASE 3
            elif mode = MinMax then paths <- paths + [MinMax(S)]     # PHASE 2
            else                    paths <- paths + [MaxMin(S)]
        T <- BruteForceConnect(paths)
        if Length(T) < best_len then
            best_len, best, hit <- Length(T), T, true
    return (best, best_len, hit)
\end{verbatim}
\end{small}

The reference implementation is \texttt{SiftTSP.py}; entry point
\texttt{sift\_tsp}.

\subsection{Prior Art and Contribution}

SiftTSP belongs to the family of \textbf{geometric-partitioning
heuristics} for Euclidean TSP, whose canonical antecedent is Karp's
1977 probabilistic analysis of partitioning algorithms~\cite{karp1977}. In that
scheme the plane is divided into rectangles, each sub-instance is
solved (optimally or heuristically), and the sub-tours are reconnected.
Later geometric heuristics, surveyed by Bentley~\cite{bentley1992}, explore strip,
space-filling-curve, and recursive variants. Arora's PTAS~\cite{arora1998} is the
theoretical apex of geometric methods, achieving a $(1+\varepsilon)$
approximation in polynomial time but with constants that make it
impractical at small $n$. State-of-the-art \emph{practical} solvers ---
LKH~\cite{helsgaun2000}, built on Lin--Kernighan~\cite{linkernighan1973}, and Concorde~\cite{applegate2006} --- are not
partitioning-based and dominate empirically on TSPLIB-scale instances.

Against this backdrop, \textbf{SiftTSP makes no theoretical claim that
extends prior art.} Its complexity bound ($O(n^2)$) is no better than
several classical heuristics; it offers no approximation ratio (unlike
Christofides~\cite{christofides1976} for metric TSP or Arora~\cite{arora1998} for Euclidean TSP); and
it is empirically not exact (\S2.3). What it offers is an
\emph{engineering recipe} combining:

\begin{enumerate}
\item a \textbf{sandclock depth schedule} $1 \to d_{\max} \to 1$ that
visits every depth twice per sweep, once on the way up and once on the
way down;
\item a \textbf{two-scale angular search} pairing a uniform ``mill''
(breadth) with a damped ``sieve'' (depth) around the mill's best
angle;
\item an \textbf{8-state Phase~4 super-cycle} that systematically
permutes section-solver mode (min-max / max-min), angular sweep
direction (CW / CCW), and visit ordering (MD / DM), with independent
doubling caps on mill and sieve that grow only on improvement.
\end{enumerate}

To the author's knowledge this specific orchestration is new, but each
individual primitive (recursive median bisection, exhaustive sub-tour
chaining, greedy section solvers, rotational axis search) is standard.
We therefore frame SiftTSP as a \textbf{deterministic, reproducible,
and falsifiable} partitioning heuristic --- and ship a concrete witness
to its non-exactness (\S2.3) rather than leave exactness as an open
question.

\subsection{Complexity}

Treat $d_{\text{ceiling}}, m_{\text{ceiling}}, s_{\text{ceiling}}, \tau$
as fixed constants independent of $n$.

\paragraph{Lemma 1 (Subdivide).} \textit{Algorithm \texttt{Subdivide}
runs in $O(n \log n)$ time and produces at most $2^d$ sections of size
$\lceil n/2^d \rceil$ or $\lfloor n/2^d \rfloor$.}

\textit{Proof.} The recursion has depth $d$. At each level the cities
in each sub-list are sorted once along the current axis and split at
the median, so each level does $O(n \log n)$ comparison work in total.
With $d$ fixed, total work is $O(n \log n)$. Median splits keep section
sizes balanced. $\square$

\paragraph{Lemma 2 (Section solving).} \textit{Let $k = 2^d$ and assume
the recursive median bisection of Lemma~1, so every section has size in
$\{\lfloor n/k \rfloor, \lceil n/k \rceil\}$. Processing all sections
takes $O(1)$ time when $\lceil n/k \rceil \leq \tau$ (every section is
brute-forced), and $O(n^2 / k)$ time otherwise; the mixed regime
(some brute-forced, some greedy) is bounded by the larger.}

\textit{Proof.} Brute force on a section of $\leq \tau$ cities is
$O(\tau!) = O(1)$, and there are at most $k = O(1)$ such sections,
giving $O(1)$ total. For the greedy branch, each section has size
$s_i \leq \lceil n/k \rceil$ by balanced bisection, and each greedy
solver (\texttt{minmax\_path}, \texttt{maxmin\_path}) does an
$O(s_i)$ scan at each of its $s_i$ steps, for $O(s_i^2)$ total per
section. Summing, $\sum_i s_i^2 \leq k \cdot (\lceil n/k \rceil)^2 =
O(n^2/k)$. $\square$

\paragraph{Lemma 3 (Connection).} \textit{For fixed $d$ (hence fixed
$k = 2^d$), \texttt{BruteForceConnect} runs in $O(n)$ time.}

\textit{Proof.} There are $(k-1)!$ orderings and $2^{k-1}$ flip vectors
(both $O(1)$ for fixed $d$), and each candidate tour is scored in
$O(n)$. $\square$

\paragraph{Lemma 4 (Phase~1).} \textit{The greedy-longest construction
in the reference implementation runs in $O(n^2)$ time.}

\textit{Proof.} At each of $n$ steps it scans the remaining unvisited
cities to find the farthest from the current endpoint and removes it
from a list, both $O(n)$ operations. $\square$

\paragraph{Theorem 1 (Time complexity of SiftTSP).} \textit{For fixed
parameters, SiftTSP runs in $O(n^2)$ time on every input.}

\textit{Proof.} Phase~1 costs $O(n^2)$ (Lemma~4). Each \texttt{EvalAt}
invocation loops over the four axis-pairs (\S1.4), each costing
$O(n \log n)$ for \texttt{Subdivide} (Lemma~1), $O(n^2/k)$ for section
solving (Lemma~2; $O(1)$ in the bounded-$n$ all-exact regime), and
$O(n)$ for \texttt{BruteForceConnect} (Lemma~3); together $O(n^2)$ per
\texttt{EvalAt}. Because $d_{\max} \leq d_{\text{ceiling}}$ throughout
the run, the number of \texttt{EvalAt} invocations per
\texttt{SandclockSweep} is
$O(d_{\text{ceiling}} \cdot m_{\text{ceiling}} + s_{\text{ceiling}}) = O(1)$.
The total number of sweeps in Phase~4 is $O(1)$ for fixed parameters
by the termination argument of \S1.3 (the per-sweep configuration is
drawn from an $O(1)$-sized set and the running best strictly
decreases along improvements).Multiplying gives $O(n^2)$.
$\square$

\paragraph{Corollary 1.} \textit{SiftTSP is a polynomial-time algorithm
for all fixed parameters.} Polynomiality is a property of a
\emph{heuristic}: \S2.3 shows it is not exact in general.

\paragraph{Remark (sharper bounds in sub-regimes).} Phase~1 dominates
the worst-case bound. If Phase~1 were replaced by an $O(n \log n)$
implementation (for example via a $k$-d tree for farthest-point
queries), and if $n$ is small enough that every section is brute-forced
($\lceil n / 2^{d_{\text{ceiling}}} \rceil \leq \tau$), the overall
complexity would drop to $O(n \log n)$. With fixed parameters that
regime is reached only for $n \leq \tau \cdot 2^{d_{\text{ceiling}}}$,
that is, $n$ bounded by a constant, in which case the entire algorithm
is trivially $O(1)$. We therefore report the honest worst-case bound
$O(n^2)$.

\section{Experiments}

\subsection{Setup}

All experiments use the Euclidean distance metric. Optimal tours are
computed by the exact Held--Karp DP~\cite{heldkarp1962} (feasible up to $n \approx 20$).
We compare SiftTSP (\texttt{SiftTSP.py}, entry point \texttt{sift\_tsp})
against:

\begin{itemize}
\item \textbf{Held--Karp}: exact $O(n^2 \cdot 2^n)$ DP --- ground-truth
optimum.
\item \textbf{Nearest-Neighbor (NN)}: classic $O(n^2)$ greedy
construction.
\item \textbf{2-opt}: 2-opt local search starting from the NN tour.
\end{itemize}

SiftTSP is run with $\tau = 6$ and $d_{\max} \in \{1, 2, 3\}$. The
optimality gap is $(L / L^{\star} - 1) \times 100\%$ where $L^{\star}$ is
the Held--Karp optimum.

Reproduce with:
\begin{small}
\begin{verbatim}
$env:PYTHONIOENCODING="utf-8"; python -m tests.adversarial   # full battery
$env:PYTHONIOENCODING="utf-8"; python counterexample.py      # canonical counterexample
\end{verbatim}
\end{small}

\paragraph{Reproducibility.} All experiments run under CPython~3.11 (no
third-party numerical libraries). The algorithm is deterministic given
$(C, d_{\text{ceiling}}, \tau, m_{\text{ceiling}}, s_{\text{ceiling}})$
and the city input ordering, with no randomised tie-breaking. Random
instance families use \texttt{random.Random} with the following fixed
seeds (\texttt{tests/adversarial.py}): \texttt{zigzag} = 0,
\texttt{interleaved} = 1, \texttt{comb} = 2, \texttt{rings} = 3,
\texttt{uniform-12/14/16} = 42/43/44. The deterministic instances
(\texttt{rings-12} of \texttt{counterexample.py}) use literal
coordinates listed in the source file.

\subsection{Adversarial Battery}

Instance families designed to probe geometric-bisection weaknesses:

\begin{itemize}
\item \textbf{zigzag}-\textit{n}: $n$ points alternating across a
vertical axis.
\item \textbf{interleaved}-\textit{n}: two clusters with interleaved
indices.
\item \textbf{comb}-\textit{n}: a spine with vertical teeth.
\item \textbf{rings}-\textit{n}: two concentric rings --- no axis is
preferred.
\item \textbf{uniform}-\textit{n}: uniform random control.
\end{itemize}

\paragraph{Table 1.} Optimality gap (\%) --- \textbf{bold = exact},
\textit{italic = sub-optimal}.

\begin{center}
\begin{tabular}{lcccccc}
\toprule
Instance & $n$ & NN & NN+2-opt & $d{=}1$ & $d{=}2$ & $d{=}3$ \\
\midrule
zigzag-12        & 12 & \textbf{0.00} & \textbf{0.00} & \textbf{0.00} & \textbf{0.00} & \textbf{0.00} \\
zigzag-14        & 14 & 28.84         & \textbf{0.00} & \textbf{0.00} & \textbf{0.00} & \textbf{0.00} \\
zigzag-16        & 16 & \textbf{0.00} & \textbf{0.00} & \textbf{0.00} & \textbf{0.00} & \textbf{0.00} \\
zigzag-tight-12  & 12 & \textbf{0.00} & \textbf{0.00} & \textbf{0.00} & \textbf{0.00} & \textbf{0.00} \\
zigzag-wide-14   & 14 & 19.52         & \textbf{0.00} & \textbf{0.00} & \textbf{0.00} & \textbf{0.00} \\
interleaved-12   & 12 & 38.09         & \textbf{0.00} & \textbf{0.00} & \textbf{0.00} & \textbf{0.00} \\
interleaved-14   & 14 & \textbf{0.00} & \textbf{0.00} & \textit{31.41} & \textbf{0.00} & \textbf{0.00} \\
comb-12          & 12 & \textbf{0.00} & \textbf{0.00} & \textbf{0.00} & \textbf{0.00} & \textbf{0.00} \\
comb-14          & 14 & \textbf{0.00} & \textbf{0.00} & \textit{25.98} & \textbf{0.00} & \textbf{0.00} \\
\textbf{rings-12} & 12 & 17.63 & 17.63 & \textit{6.61} & \textit{6.61} & \textit{3.24} \\
\textbf{rings-14} & 14 & 8.58  & 8.58  & \textit{26.06} & \textit{3.82} & \textit{1.52} \\
uniform-12       & 12 & 12.20 & \textbf{0.00} & \textbf{0.00} & \textit{5.46} & \textbf{0.00} \\
uniform-14       & 14 & 30.36 & \textbf{0.00} & \textit{15.32} & \textit{1.47} & \textit{1.47} \\
uniform-16       & 16 & 12.01 & 1.01          & \textit{9.78}  & \textit{7.60} & \textit{0.31} \\
\bottomrule
\end{tabular}
\end{center}

\subsection{Counterexample to Exactness --- \texttt{rings-12}}

We exhibit an explicit 12-point Euclidean instance on which SiftTSP
fails to recover the optimum at every tested depth. The instance, an
inner ring of 6 points of radius $\approx 2$ and an outer ring of 6
points of radius $\approx 4$ both centred at $(5,5)$, is materialised
verbatim in \texttt{counterexample.py}.

\paragraph{Table 2.} Counterexample \texttt{rings-12} --- Held--Karp
optimum $L^{\star} = 29.7823$.

\begin{center}
\begin{tabular}{lcc}
\toprule
Algorithm & Length $L$ & Gap \\
\midrule
Held--Karp (exact)            & 29.7823 & --- \\
SiftTSP, $d=1$, $\tau=6$      & 31.7503 & $+6.61\%$ \\
SiftTSP, $d=2$, $\tau=6$      & 31.7503 & $+6.61\%$ \\
SiftTSP, $d=3$, $\tau=6$      & 30.7478 & $+3.24\%$ \\
\bottomrule
\end{tabular}
\end{center}

\paragraph{Proposition 1 (Non-exactness).} \textit{SiftTSP with
$\tau = 6$ and any $d \in \{1, 2, 3\}$ is not exact on the Euclidean
TSP.}

\textit{Proof.} The 12-point instance \texttt{RINGS\_12} in
\texttt{counterexample.py} is a witness: Held--Karp's optimum has
length $29.7823$, while SiftTSP at $d=1,2$ returns $31.7503$ and at
$d=3$ returns $30.7478$, all strictly greater. Values are reproducible
by \texttt{python counterexample.py}. $\square$

The following is the structural reason \emph{why} \texttt{rings-12}
defeats the algorithm. Proposition~1 above is empirical; Proposition~2
below is a structural impossibility result that applies to a broad
class of algorithms including SiftTSP.

\paragraph{Proposition 2 (Decomposition irrecoverability).}
\textit{Let $C \subset \mathbb{R}^2$ be a finite point set and let
$\mathcal{D}$ be a finite family of partitions of $C$, each
$D \in \mathcal{D}$ a partition $C = S_1^D \sqcup \dots \sqcup S_{k_D}^D$
into pairwise disjoint components. Consider any algorithm $\mathcal{A}$
that, for each $D \in \mathcal{D}$, computes an open Hamiltonian path
on each component of $D$ and then closes those paths into a tour by
chaining them in some order with some orientation per path, and
finally returns the shortest such candidate over all $D \in \mathcal{D}$
and all chainings. Call a tour $\pi$ on $C$ \textbf{component-contiguous}
for $D$ if, for every component $S_i^D$, the cities of $S_i^D$ occupy a
contiguous arc of $\pi$. If for every $D \in \mathcal{D}$ the global
optimum $\pi^{\star}$ fails to be component-contiguous for $D$, then
$L(\mathcal{A}(C)) > L(\pi^{\star})$.}

\textit{Proof.} For a fixed $D \in \mathcal{D}$ every candidate tour
considered by $\mathcal{A}$ visits each component on a single
contiguous arc, by construction (the open Hamiltonian path of a
component is contiguous, and chaining concatenates components without
interleaving them). Hence every candidate is component-contiguous for
$D$. By hypothesis $\pi^{\star}$ is not component-contiguous for any
$D \in \mathcal{D}$, so $\pi^{\star}$ is never among the candidates,
and $\mathcal{A}(C)$ --- the minimum over the candidates --- has length
strictly greater than $L(\pi^{\star})$ unless some non-optimal candidate
happens to have length exactly $L(\pi^{\star})$. In the Euclidean
generic position (no two distinct tours have equal length) this
exception is empty. $\square$

Proposition~2 applies to SiftTSP because the family $\mathcal{D}$ it
explores (median bisections at depths $1, \dots, d_{\max}$ along
finitely many rotated axes and four axis-pairs) is \emph{fixed in
advance} and does not depend on $\pi^{\star}$. The \texttt{rings-12}
Held--Karp optimum is empirically non-contiguous with respect to every
$D$ tried at $d \leq 3$: any straight-line bisection either separates
the two rings or cuts each ring into two diametrically opposed
half-rings, while $\pi^{\star}$ alternates between inner and outer
rings. Finer angular search adds resolution but does not change the
decomposition primitive, so it cannot defeat the obstruction.

\subsection{Observations}

The following are observations from the 14-instance battery (Table~1).
Each instance is run on a single fixed seed; we report descriptive
findings, not statistical claims.

\begin{enumerate}
\item \textbf{Structured inputs are matched.} On zigzag, comb, and
interleaved instances ($n \in \{12, 14, 16\}$), SiftTSP at $d \geq 2$
matches Held--Karp in every tested case, including instances where NN
has gaps up to $+38\%$.
\item \textbf{Two rotationally symmetric inputs are not matched, but
2-opt is beaten on both.} On \texttt{rings-12} and \texttt{rings-14},
NN and 2-opt are stuck at $+17.63\%$ and $+8.58\%$; SiftTSP at $d=3$
achieves $+3.24\%$ and $+1.52\%$. Sample size is two; we do not claim
this pattern holds on a wider population of rotationally symmetric
instances.
\item \textbf{Increasing depth is not monotone.} On \texttt{uniform-12},
$d{=}2$ gives $+5.46\%$ while $d \in \{1, 3\}$ are exact. Finer
subdivision can split clusters that the optimum traverses contiguously.
\end{enumerate}

\subsection{Discussion}

\paragraph{Where SiftTSP appears useful (on this battery).} Determinism
(reproducible given $(C, d, \tau, m, s)$); strong performance on inputs
with a clear geometric scan order (zigzag/comb/interleaved); and a
suggestive regime on the two rotationally symmetric instances tested,
where 2-opt is trapped. A larger study would be needed to claim this
generalises.

\paragraph{Limits.} Decomposition irrecoverability (\S2.3) is the
structural ceiling of \emph{any} algorithm that explores only a
fixed-in-advance family of geometric decompositions and then optimally
chains the resulting component paths. The \texttt{rings-12} $+3.24\%$
wall is unmoved by the mill+sieve refinement, consistent with the
obstruction being structural rather than angular-resolution-limited.
The connection step's $(2^d - 1)! \cdot 2^{2^d - 1}$ growth blocks
$d \geq 4$ unless replaced by a smarter chaining primitive.

\paragraph{Open directions.}
\begin{itemize}
\item \emph{Adaptive decomposition} (axes that depend on local
geometry).
\item \emph{Recursive connection} (replace the brute-force chain by a
recursive SiftTSP application on the path endpoints).
\item \emph{Hybridisation} with 2-opt or Lin--Kernighan~\cite{linkernighan1973}
(LKH~\cite{helsgaun2000}) using SiftTSP as an initial-tour generator.
\item \emph{Characterisation of exactness regimes} --- a sufficient
geometric condition on $C$ guaranteeing exactness would clarify the
value of the algorithm as a conditional exact solver.
\end{itemize}

None of these are claimed to defeat the decomposition-irrecoverability
obstruction. SiftTSP is and will remain a heuristic; this work makes
\textbf{no claim} of resolving $\mathcal{P}$ versus $\mathcal{NP}$.

\section*{Code and Data Availability}

The reference implementation, all instance families, and the
verification scripts are available at
\url{https://github.com/McNopper/SiftTSP}. Specifically:
\texttt{SiftTSP.py} (algorithm), \texttt{tests/adversarial.py} (Table~1
battery, fixed seeds as above), \texttt{counterexample.py} (Table~2
\texttt{rings-12} instance with literal coordinates and Held--Karp
verification). Numbers reported in this paper were produced by these
scripts under CPython~3.11.

\section*{References}

\begin{thebibliography}{99}\small

\bibitem{cook1971}
S.~A.~Cook, ``The complexity of theorem-proving procedures,''
in \emph{Proceedings of the 3rd Annual ACM Symposium on Theory of
Computing (STOC)}, pp.~151--158. ACM, New York (1971).
\url{https://doi.org/10.1145/800157.805047}

\bibitem{levin1973}
L.~A.~Levin, ``Universal sequential search problems,''
\emph{Problems of Information Transmission} \textbf{9}(3), 265--266 (1973).

\bibitem{aaronson2009}
S.~Aaronson, ``P vs.\ NP,''
in R.~Downey (ed.) \emph{Fundamentals of Computation Theory}.
Springer, Berlin (2009).
\url{https://www.scottaaronson.com/papers/pnp.pdf}

\bibitem{applegate2006}
D.~L.~Applegate, R.~E.~Bixby, V.~Chv\'atal, and W.~J.~Cook,
\emph{The Traveling Salesman Problem: A Computational Study}.
Princeton University Press, Princeton (2006).

\bibitem{karp1972}
R.~M.~Karp, ``Reducibility among combinatorial problems,''
in R.~E.~Miller and J.~W.~Thatcher (eds.) \emph{Complexity of Computer
Computations}, pp.~85--103. Plenum Press, New York (1972).
\url{https://doi.org/10.1007/978-1-4684-2001-2_9}

\bibitem{papadimitriou1977}
C.~H.~Papadimitriou, ``The Euclidean travelling salesman problem is
NP-complete,''
\emph{Theoretical Computer Science} \textbf{4}(3), 237--244 (1977).
\url{https://doi.org/10.1016/0304-3975(77)90012-3}

\bibitem{heldkarp1962}
M.~Held and R.~M.~Karp, ``A dynamic programming approach to
sequencing problems,''
\emph{Journal of the Society for Industrial and Applied Mathematics}
\textbf{10}(1), 196--210 (1962).
\url{https://doi.org/10.1137/0110015}

\bibitem{karp1977}
R.~M.~Karp, ``Probabilistic analysis of partitioning algorithms for
the traveling-salesman problem in the plane,''
\emph{Mathematics of Operations Research} \textbf{2}(3), 209--224 (1977).
\url{https://doi.org/10.1287/moor.2.3.209}

\bibitem{christofides1976}
N.~Christofides, ``Worst-case analysis of a new heuristic for the
travelling salesman problem,''
Report 388, Graduate School of Industrial Administration,
Carnegie Mellon University (1976).

\bibitem{linkernighan1973}
S.~Lin and B.~W.~Kernighan, ``An effective heuristic algorithm for
the traveling-salesman problem,''
\emph{Operations Research} \textbf{21}(2), 498--516 (1973).
\url{https://doi.org/10.1287/opre.21.2.498}

\bibitem{helsgaun2000}
K.~Helsgaun, ``An effective implementation of the Lin--Kernighan
traveling salesman heuristic,''
\emph{European Journal of Operational Research} \textbf{126}(1),
106--130 (2000).
\url{https://doi.org/10.1016/S0377-2217(99)00284-2}

\bibitem{arora1998}
S.~Arora, ``Polynomial time approximation schemes for Euclidean
traveling salesman and other geometric problems,''
\emph{Journal of the ACM} \textbf{45}(5), 753--782 (1998).
\url{https://doi.org/10.1145/290179.290180}

\bibitem{bentley1992}
J.~L.~Bentley, ``Fast algorithms for geometric traveling salesman
problems,''
\emph{ORSA Journal on Computing} \textbf{4}(4), 387--411 (1992).
\url{https://doi.org/10.1287/ijoc.4.4.387}

\end{thebibliography}

\end{document}
